\documentclass[11pt]{article}
% ===========================================================================
%  Shared preamble for all MPM2D worksheets — input right after \documentclass.
%  (Files beginning with "_" are skipped by mpm2d-worksheet-publish.mjs.)
% ===========================================================================
\usepackage[letterpaper,top=2.2cm,bottom=1.9cm,left=1.6cm,right=1.6cm,headheight=28pt,footskip=24pt]{geometry}
\usepackage{amsmath,amssymb}
\usepackage[T1]{fontenc}
\usepackage[utf8]{inputenc}
\usepackage{lmodern}
\usepackage{xcolor}
\usepackage[most]{tcolorbox}
\usepackage{enumitem}
\usepackage{multicol}
\usepackage{fancyhdr}
\usepackage{tikz}
\usetikzlibrary{calc}
\usepackage{pgfplots}
\pgfplotsset{compat=1.18}

\definecolor{exblue}{HTML}{4A90E2}
\definecolor{exbg}{HTML}{E6F3FF}
\definecolor{qorange}{HTML}{E69138}
\definecolor{qbg}{HTML}{FFF7CC}
\definecolor{keygray}{HTML}{475569}
\definecolor{keybg}{HTML}{F1F5F9}
\definecolor{iagreen}{HTML}{14653B}
\definecolor{titleblue}{HTML}{1D4ED8}
% Rotating palette so each worked example gets its own colour.
\definecolor{ex0}{HTML}{2563A0}\definecolor{ex1}{HTML}{3B7D3B}\definecolor{ex2}{HTML}{8A5A00}
\definecolor{ex3}{HTML}{A3327A}\definecolor{ex4}{HTML}{6D28A3}\definecolor{ex5}{HTML}{B08900}
\definecolor{ex6}{HTML}{0E7490}\definecolor{ex7}{HTML}{9A3412}\definecolor{ex8}{HTML}{15803D}

\pagestyle{fancy}
\fancyhf{}
\fancyhead[L]{\footnotesize NAME: \underline{\hspace{3.4cm}}}
\fancyhead[C]{\textbf{Integration Academy}\\[-2pt]{\scriptsize Math Simplified \textperiodcentered{} Success Amplified}}
\fancyhead[R]{\footnotesize MPM2D}
\fancyfoot[L]{\scriptsize dr.merie@IntegrationAcademy.ca}
\fancyfoot[C]{\scriptsize\thepage}
\fancyfoot[R]{\scriptsize IntegrationAcademy.ca}
\renewcommand{\headrulewidth}{1.2pt}
\renewcommand{\footrulewidth}{0.4pt}
\setlength{\parindent}{0pt}
\setlength{\parskip}{4pt}

% Examples are NOT breakable, so each example + its solution stays on one page.
\newtcolorbox{example}[2]{colframe=#1,colback=#1!7,coltitle=white,colbacktitle=#1,fonttitle=\bfseries,title={#2},arc=3pt,boxrule=1pt}
\newtcolorbox{qbox}[1]{colframe=qorange,colback=qbg,coltitle=white,colbacktitle=qorange,fonttitle=\bfseries,title={#1},arc=3pt,breakable,boxrule=1pt}
\newtcolorbox{keybox}{colframe=keygray,colback=keybg,coltitle=white,colbacktitle=keygray,fonttitle=\bfseries,title={$\checkmark$\ Answer Key},arc=3pt,breakable,boxrule=1pt}
\newtcolorbox{recapbox}{colframe=keygray,colback=keybg,coltitle=white,colbacktitle=keygray,fonttitle=\bfseries,title={Question Recap},arc=3pt,breakable,boxrule=1pt}
\newtcolorbox{ideabox}{colframe=titleblue,colback=titleblue!6,coltitle=white,colbacktitle=titleblue,fonttitle=\bfseries,title={The Big Ideas},arc=3pt,breakable,boxrule=1.2pt}
\newtcolorbox{learnbox}{colframe=iagreen,colback=iagreen!5,coltitle=white,colbacktitle=iagreen,fonttitle=\bfseries,title={Core Concepts},arc=3pt,breakable,boxrule=1.2pt}
\newcommand{\lh}[1]{\par\smallskip{\bfseries\color{iagreen}#1}\par\nobreak\smallskip}

\newcommand{\soln}{\par\smallskip\textit{\textbf{Solution.}}\ }
\newcommand{\workspace}[1]{\par\vspace{3pt}\fcolorbox{black!30}{white}{\begin{minipage}[t][#1][t]{\dimexpr\linewidth-2\fboxsep\relax}\ \end{minipage}}\par}
\newcommand{\sech}[1]{\par\vspace{8pt}{\Large\bfseries\color{iagreen}#1}\par\nobreak\vspace{2pt}{\color{black!20}\hrule height1pt}\vspace{6pt}}

% A compact coordinate plot sized for an example box: \eplot{xmin}{xmax}{ymin}{ymax}{body}
\newcommand{\eplot}[5]{\begin{center}\begin{tikzpicture}
\begin{axis}[width=6.8cm,height=4.4cm,axis lines=middle,xlabel={\tiny$x$},ylabel={\tiny$y$},
grid=both,grid style={gray!16},xmin=#1,xmax=#2,ymin=#3,ymax=#4,
xtick distance=2,ytick distance=2,tick label style={font=\tiny},axis line style={-},samples=120]
#5
\end{axis}\end{tikzpicture}\end{center}}

% Same as \eplot but with its own tick spacing: \eplotx{xmin}{xmax}{ymin}{ymax}{xtick}{ytick}{body}
\newcommand{\eplotx}[7]{\begin{center}\begin{tikzpicture}
\begin{axis}[width=8.4cm,height=5.4cm,axis lines=middle,xlabel={\tiny$x$},ylabel={\tiny$y$},
grid=both,grid style={gray!16},xmin=#1,xmax=#2,ymin=#3,ymax=#4,
xtick distance=#5,ytick distance=#6,tick label style={font=\tiny},axis line style={-},samples=120]
#7
\end{axis}\end{tikzpicture}\end{center}}

% A sine-curve plot (degrees on x, 0..360): \splot{ymin}{ymax}{body}
\newcommand{\splot}[3]{\begin{center}\begin{tikzpicture}
\begin{axis}[width=8.6cm,height=4.6cm,axis lines=middle,xlabel={\tiny$x\,(^\circ)$},ylabel={\tiny$y$},
grid=both,grid style={gray!16},xmin=0,xmax=360,ymin=#1,ymax=#2,
xtick distance=90,ytick distance=1,tick label style={font=\tiny},axis line style={-},samples=180]
#3
\end{axis}\end{tikzpicture}\end{center}}

% A small coordinate plot: \plot{xmin}{xmax}{ymin}{ymax}{body}
\newcommand{\plot}[5]{\begin{center}\begin{tikzpicture}
\begin{axis}[width=8.4cm,height=6cm,axis lines=middle,xlabel={\scriptsize$x$},ylabel={\scriptsize$y$},
grid=both,grid style={gray!16},xmin=#1,xmax=#2,ymin=#3,ymax=#4,
xtick distance=2,ytick distance=2,tick label style={font=\tiny},axis line style={-}]
#5
\end{axis}\end{tikzpicture}\end{center}}

% Right triangle (angle theta at bottom-left, right angle at bottom-right):
%   \rtri{drawBase}{drawHeight}{oppLabel}{adjLabel}{hypLabel}
\newcommand{\rtri}[5]{\begin{center}\begin{tikzpicture}[scale=0.85]
\coordinate (A) at (0,0);\coordinate (B) at (#1,0);\coordinate (C) at (#1,#2);
\draw[thick,exblue] (A)--(B)--(C)--cycle;
\draw[black] ($(B)+(-0.32,0)$)--($(B)+(-0.32,0.32)$)--($(B)+(0,0.32)$);
\node[below] at ($(A)!0.5!(B)$){#4};
\node[right] at ($(B)!0.5!(C)$){#3};
\node[above left] at ($(A)!0.5!(C)$){#5};
\node at (0.7,0.22){$\theta$};
\end{tikzpicture}\end{center}}

% General (oblique) triangle with vertex labels and opposite-side labels:
%   \gtri{Alabel}{Blabel}{Clabel}{aSide}{bSide}{cSide}
\newcommand{\gtri}[6]{\begin{center}\begin{tikzpicture}[scale=0.85]
\coordinate (A) at (0,0);\coordinate (B) at (5,0);\coordinate (C) at (1.6,3);
\draw[thick,exblue] (A)--(B)--(C)--cycle;
\node[below left] at (A){#1};\node[below right] at (B){#2};\node[above] at (C){#3};
\node[below] at ($(A)!0.5!(B)$){#6};
\node[right] at ($(B)!0.5!(C)$){#4};
\node[left] at ($(A)!0.5!(C)$){#5};
\end{tikzpicture}\end{center}}

\fancyhead[R]{\footnotesize MHF4U \textperiodcentered{} 7.1}
\begin{document}

\begin{center}
{\LARGE\bfseries\color{titleblue}Worksheet 7.1 --- Average \& Instantaneous Rate of Change}\\[2pt]
{\small Grade 12 Advanced Functions (MHF4U) \textperiodcentered{} Unit 7: Rates of Change and Combining Functions}
\end{center}

\begin{ideabox}
The average rate of change over $[a,b]$ is the secant slope $\dfrac{f(b)-f(a)}{b-a}$; the instantaneous rate at a point is the tangent slope, estimated over a tiny interval.
\begin{itemize}[leftmargin=*,itemsep=2pt]
  \item Average rate $=\dfrac{f(b)-f(a)}{b-a}$ (secant slope).
  \item Instantaneous rate $=$ tangent slope at a point.
  \item Estimate it with a small interval like $[a,a+0.1]$.
\end{itemize}
\end{ideabox}

\begin{learnbox}
\lh{Average rate of change}The \textbf{average rate of change} over an interval is the slope of the secant line, $\dfrac{f(b)-f(a)}{b-a}$.
\lh{Instantaneous rate}The \textbf{instantaneous rate} at a point is the slope of the tangent line, approached by shrinking the interval toward that point.
\lh{Estimating it}Use average rates over smaller and smaller intervals around a point to estimate the instantaneous rate there.
\end{learnbox}

\sech{Worked Examples}

\begin{example}{ex0}{Example 1 --- Average rate}
Average rate of change of $f(x)=x^2-4$ on $[0,3]$.\soln $\dfrac{f(3)-f(0)}{3-0}=\dfrac{5-(-4)}{3}=3$. The secant slope approaches the tangent as the interval shrinks:\eplot{-3}{4}{-5}{7}{\addplot[exblue,very thick,domain=-2.8:3.2]{x^2-4};\addplot[red,thick] coordinates {(0,-4) (3,5)};\addplot[red,only marks,mark=*,mark size=1.6pt] coordinates {(0,-4) (3,5)};}
\end{example}

\begin{example}{ex1}{Example 2 --- Average rate}
Average rate of change of $f(x)=3x^2$ on $[1,3]$.\soln $\dfrac{27-3}{2}=12$.
\end{example}

\begin{example}{ex2}{Example 3 --- Estimate instantaneous}
Estimate the instantaneous rate of $f(x)=x^2$ at $x=4$ using $[4,4.1]$.\soln $\dfrac{4.1^2-4^2}{0.1}=\dfrac{16.81-16}{0.1}=8.1$ (true value $8$).
\end{example}

\begin{example}{ex3}{Example 4 --- Linear rate}
Average rate of change of $f(x)=4x-3$ on $[1,6]$.\soln $\dfrac{21-1}{5}=4$ — constant for a line.
\end{example}

\begin{example}{ex4}{Example 5 --- Interpret}
A car's position is $d(t)$. What is $\dfrac{d(9)-d(3)}{9-3}$?\soln The average speed over $[3,9]$.
\end{example}

\begin{example}{ex5}{Example 6 --- Average rate}
Average rate of $f(x)=x^2$ on $[3,7]$.\soln $\dfrac{49-9}{4}=10$.
\end{example}

\begin{example}{ex6}{Example 7 --- Cubic}
Average rate of $f(x)=x^3$ on $[2,3]$.\soln $\dfrac{27-8}{1}=19$.
\end{example}

\begin{example}{ex7}{Example 8 --- Estimate instantaneous}
Estimate the instantaneous rate of $f(x)=x^3$ at $x=1$ using $[1,1.1]$.\soln $\dfrac{1.331-1}{0.1}=3.31\approx3$.
\end{example}

\begin{example}{ex8}{Example 9 --- Tangent meaning}
What does a tangent line's slope measure?\soln The instantaneous rate of change at that point.
\end{example}

\clearpage
\sech{Your Turn --- Show Your Work}
For each question, show all steps clearly in the space provided.

\begin{qbox}{Question 1}
Average rate of $f(x)=x^2-4x$ on $[1,5]$?\workspace{2.4cm}
\end{qbox}

\begin{qbox}{Question 2}
Average rate of $f(x)=2x^3$ on $[1,2]$?\workspace{2.4cm}
\end{qbox}

\begin{qbox}{Question 3}
Estimate the instantaneous rate of $f(x)=x^2$ at $x=5$ using $[5,5.1]$.\workspace{2.4cm}
\end{qbox}

\begin{qbox}{Question 4}
Average rate of $f(x)=7x+1$ on $[2,5]$?\workspace{2.4cm}
\end{qbox}

\begin{qbox}{Question 5}
What does the slope of a tangent line measure?\workspace{2.4cm}
\end{qbox}

\begin{qbox}{Question 6}
Average rate of $f(x)=x^2$ on $[2,8]$?\workspace{2.4cm}
\end{qbox}

\begin{qbox}{Question 7}
Average rate of $f(x)=x^2+1$ on $[0,3]$?\workspace{2.4cm}
\end{qbox}

\begin{qbox}{Question 8}
Average rate of $f(x)=3x$ on $[2,6]$?\workspace{2.4cm}
\end{qbox}

\begin{qbox}{Question 9}
Estimate the instantaneous rate of $f(x)=x^2$ at $x=6$ using $[6,6.1]$.\workspace{2.4cm}
\end{qbox}

\begin{qbox}{Question 10}
Average rate of $f(x)=x^3$ on $[0,2]$?\workspace{2.4cm}
\end{qbox}

\begin{qbox}{Question 11}
Does a line have a changing instantaneous rate?\workspace{2.4cm}
\end{qbox}

\begin{qbox}{Question 12}
Average rate of $f(x)=-x^2$ on $[1,3]$?\workspace{2.4cm}
\end{qbox}

\begin{qbox}{Question 13 --- Challenge}
A population is $P(t)=100+t^2$. Find its average rate of growth on $[0,5]$.\workspace{3cm}
\end{qbox}

\clearpage
\sech{Question Recap \& Answer Key}

\begin{recapbox}
\begin{multicols}{2}
\small
\begin{enumerate}[leftmargin=*,itemsep=3pt]
  \item Average rate of $f(x)=x^2-4x$ on $[1,5]$?
  \item Average rate of $f(x)=2x^3$ on $[1,2]$?
  \item Estimate the instantaneous rate of $f(x)=x^2$ at $x=5$ using $[5,5.1]$.
  \item Average rate of $f(x)=7x+1$ on $[2,5]$?
  \item What does the slope of a tangent line measure?
  \item Average rate of $f(x)=x^2$ on $[2,8]$?
  \item Average rate of $f(x)=x^2+1$ on $[0,3]$?
  \item Average rate of $f(x)=3x$ on $[2,6]$?
  \item Estimate the instantaneous rate of $f(x)=x^2$ at $x=6$ using $[6,6.1]$.
  \item Average rate of $f(x)=x^3$ on $[0,2]$?
  \item Does a line have a changing instantaneous rate?
  \item Average rate of $f(x)=-x^2$ on $[1,3]$?
  \item A population is $P(t)=100+t^2$. Find its average rate of growth on $[0,5]$.
\end{enumerate}
\end{multicols}
\end{recapbox}

\begin{keybox}
\begin{multicols}{2}
\small
\begin{enumerate}[leftmargin=*,itemsep=3pt]
  \item $2$
  \item $14$
  \item $10.1\approx10$
  \item $7$
  \item instantaneous rate of change
  \item $10$
  \item $3$
  \item $3$
  \item $12.1\approx12$
  \item $4$
  \item no (constant)
  \item $-4$
  \item $\dfrac{125-100}{5}=5$
\end{enumerate}
\end{multicols}
\end{keybox}

\end{document}
